\documentclass{jsarticle}
\usepackage[dvipdfmx, final]{graphicx}
\usepackage{amsmath}
\usepackage{amssymb}
\begin{document}
\title{Beta function reveal with global differential manifold.}
\author{Masaaki Yamaguchi}
\date{}
\maketitle
 
  Global differential manifold exclude with constant of value of imaginary
  and real to Euler law equation, and this equation equal with beta function.
  $$ {d \over {d f}}F(x,y) = {d \over {d f}}\int\int{1 \over {(x\log x)}^{2}}
  dx_m + {d \over {d f}}\int\int{1 \over {(y \log y)}^{1 \over 2}}dy_m $$
  This equation is hyper circle function. And Jones manifold.
  $$ = {1 \over 2}i\times 1 \times \sin(90^{\circ}) + {1 \over 2}\times 1\times
  1\times \sin(90^{\circ}) = \int {1 \over {\sin x}}dx_m $$
  $$ = \log(\sin x) = e^{x\log x} + e^{-x\log x} \ge e^{x\log x} - e^{-x \log x}
  = \cosh^{-1}(h) + \sinh^{-1}(h) $$
  This equation equal with beta function.
  $$ = \beta(p,q) $$
  Beta function escourt with gravity and anti-gravity equation.

  And, this equation system call function to deprivate of global manifold.
  Moreover, this system also recreate with integral manifold of global topology.

  This equation of global deprivate and integral manifold are computed being
  back explain to improve equation system with being existed from a accident
  of being nutral pond.

  $$ \int {1 \over {\sin x}}dx_m = \cos x \log(\sin x) = \log{(\sin x)}^{{\sin 
  x}^{'}} = {d \over {d f}}F(x) = F^{f^{'}} $$
  this result equation is non-certain system concerned with particle, electric 
  particle and atoms in open set group with possiblity of surface values,
  and have abled to proof in this group of system. Universe and the other
  dimension of concerned of relationship values.
  $$ {d \over df}F = {d \over df} \int \int{1 \over (x \log x)^2}dx_m 
  + {d \over df} \int \int {1 \over (y \log y)^{1 \over 2}}dy_m
  = \Gamma(x,y) $$
  $$ {d \over d \gamma}\Gamma = (e^{-x}x^{1-t})^{\gamma^{'}} \to
  {d \over d \gamma}\Gamma = \Gamma^{(\gamma)^{'}}$$



    $$ \Gamma(s) = \int e^{-x} x^{s - 1}dx, \Gamma^{'}(s)
    = \int e^{-x} x^{s - 1}\log x dx $$
    $$ x^{s - 1} = e^{(s - 1)\log x} $$
    $$ {\partial \over \partial s}x^{s - 1} =
       {\partial \over \partial s}e^{-x}\cdot (e^{{s - 1}\log x}) 
       = e^{-x}\cdot \log x \cdot e^{(s - 1) \log x} $$
    $$ = e^{-x}\cdot \log x \cdot x^{s - 1} $$
    $$ {d \over d \gamma}\Gamma = \Gamma^{(\gamma)^{'}} $$
    $$ {d \over d \gamma}\Gamma(s) = \left( \int^{\infty}_{0}e^{-x}
    x^{s - 1}dx \right)^{\left(\int^{\infty}_{0}e^{-x}x^{s - 1}\log x dx
    \right)^{'}} $$
    $$ = \Gamma^{(\Gamma \int \log x dx)^{'}} $$
    $$ = e^{-x \log x} $$

    $$ \Gamma(s) = \int^{\infty}_{0} e^{-x}x^{s - 1}dx $$
    $$  \Gamma^{'}(s) = \int^{\infty}_{0} e^{-x}x^{s - 1} \log x dx $$
    $$ {d \over df}F = \int x^{s - 1}dx $$
    $$ \int F dx_m = \int e^{-x}dx $$
    $$ {d \over df}F = F^{(f)^{'}}, \int F dx_m = F^{(f)} $$
    Global differential manifold and global integrate manifold are
    reached with begin estimate of gamma and beta function leaded solved,
    Global differential manifold is emerged with gamma function of
    themselves of first value of global differential manifold of gamma
    function, this equation of quota in newton and dalarnverles of equation
    is also of same results equations. Zeta function also resulted with
    quantum group reach with quanto algebra equation.
        Included of diverse of gravity of influent in quantum physics,
    also quantum physics doesn't exist Gauss surface of theorem,
   
    $$ H \Psi = \bigoplus (i \hbar ^{\nabla})^{\oplus L} $$
    $$ = \bigoplus {{ H \Psi} \over \nabla L} $$
    $$ = e^{x \log x} = x^{(x)^{'}} $$
    however,
    $$ {d \over df}F = m(x) $$
    and gravity equation in being abled to existed with quantum physics,
    and this physics is also represented with quantum level of differential
    geometry.
    $$ {d \over dt}\psi(t) = \hbar $$
    $$ = {1 \over 2}i e^{i \Hat{H}} $$
    $$ (i \hbar)^{'} = (- e^{i \Hat{H}})^{'} $$
    $$ = - i e^{i \Hat{H}} $$
    $$ \psi(x) = e^{-i \Hat{H}t}, \bigoplus (i \hbar ^{\nabla})^{\oplus L} 
    = {{{1 \over 2} e^{i \Hat{H}}}^{(-i e^{i \Hat{H}})}} $$
    $$ = ({1 \over 2}f)^{-i f}$$
    $$ = ({1 \over 2})^{-i f} \cdot e^{-x \log x} \cdot (f)^{i} $$
    $$ = \int e^{- x} x^{t - 1}dx, {d \over d \gamma}\Gamma = e^{-x \log x}  $$
     
    $$ f = x, i = t, {1 \over 2} = a, \bigoplus a^{-t x}x^t [I_m]
     \cong \int e^{-x} x^{t - 1}dx $$
     Quantum level of differential geomerty is also resulted with
    Gamma function of global differential manifold of values, Heisenberg
    equation is alos resulted with global integrate manifold and this
    quantum level of differential geometry is manifold of deprive of formula.
    $$ |\psi(t)\rangle_s = e^{-i\Hat{H} t}|\Psi\rangle_{H},
    \Hat{A}_s = \Hat{A}_H(0) $$
    $$ |\Psi(t)\rangle_s \to {d \over dt} $$
    $$ i {d \over dt}|\psi (t) \rangle_s = \Hat{H}|\psi(t)\rangle_s $$
    $$ \langle\Hat{A}(t)\rangle = \langle\Psi(t)|\Hat{A}(0)|\Psi(t)\rangle  $$
    $$ {d \over dt}\Hat{A} = {1 \over i}[\Hat{A},H] $$
    $$ \Hat{A}(t) = e^{i\Hat{H}t}\Hat{A}(0)e^{-i\Hat{H}t} $$
    $$ \lim_{\theta \to 0} \begin{pmatrix} \sin \theta \\
    \cos \theta \end{pmatrix} \begin{pmatrix} \theta & 1 \\
    1 & \theta \end{pmatrix} \begin{pmatrix} \cos \theta \\
    \sin \theta \end{pmatrix} = \begin{pmatrix} 1 & 0 \\
    0 & - 1 \end{pmatrix} $$
    $$ f^{-1}(x) x f(x) = I^{'}_m, I^{'}_m = [1,0] \times [0,1] $$

    $$ {x + y} \ge \sqrt{x y} $$
    $$ {x^{{1 \over 2} + iy} \over e^{x \log x}} = 1 $$
    $$ \mathcal{O}(x) = \nabla_i \nabla_j \int e^{{2 \over m}\sin \theta 
    \cos \theta } \times {{N \mathrm{mod}(e^{x \log x})} \over {O(x)(x +
    \Delta |f|^2)^{1 \over 2}}} $$
    $$ x \Gamma(x) = 2 \int|\sin 2 \theta|^2 d \theta, \mathcal{O}(x)
    = m(x)[D^2 \psi] $$
    $$ i^2 = (0,1) \cdot (0,1), |a||b|\cos \theta = - 1 $$
    $$ E = \mathrm{div}(E, E_1) $$
    $$ \left({\{f,g\}} \over [f,g] \right) = i^2, E = mc^2, I^{'} = i^2 $$
     Gamma function of global differential equation is first of differential
    function deprivate with ordinary differential equations, more also
    universe and other dimensiion of Higgs field emerge with zeta function,
    more spectrum focus is three dimension of manifold of singularity or pair
    theorem, and universe and other dimension of each value equals.
    This reverse of circle function of manifolds is also Higgs fields of 
    dense of entropy and result equation. And differential structure of quantum
    level in gravity is poled with zeta function and quantum group in high 
    result values. Zeta funtion and quantum group is Global differential
    manifold of result in values.
    $$ {d \over d \gamma} \Gamma = m(x) $$
    $$ = e^{- x \log x} $$
    $$ \sin i x = {{e^{-x} + e^{x}} \over 2i} $$
    $$ {d \over df}F = m(x) = e^{f} + e^{-f} $$
    $$ = 2 i \sin (i {x \log x}) $$

    Hyper circle function is constructed with sheap of manifold in 
    Beta function of reverse system emerged from Gamma fucntion 
    and reverse of circle function of entropy value,
    and this reverse of hyper circle function of beta system
    is universe and the other dimension of zeta function,
    more over spectrum focus is this beta function of reverse in
    circle function is quantum level of differential structure oneselves,
    and this quantum level is also gamma function themselves.
    This gamma fucntion is more also D-brane and anti-D-brane built 
    with supplicant values. Higgs fileds is more also pair of universe and
    other dimension each other value elements. 


    My son in acasicrecord demonstrate with being have with quantum level
    of differential structure equation.
    I have with same equation already from being on 
    my father and my doctors of professors give me hints with 
    this qlds equation explanade with global differential manifold.

    $$ H\Psi = \bigoplus {(i \hbar^{\nabla})}^{\oplus L}  (1)$$
    $$ = i\hbar\psi $$
    $$ {d \over df}F = m(x) (2) $$
    $$ = {d \over df}\int \int{1 \over (x \log x)^2}dx_m
    + {d \over df}\int \int{1 \over (y \log y)^{1 \over 2}}dy_m  (3)$$
    $$  \bigoplus {(i \hbar^{\nabla})}^{\oplus L} = \int e^{-x}x^{1 - t}dx_m 
    (4)$$
    $$ = {d \over d \gamma}\Gamma = \int \Gamma(\gamma)^{'}dx_m (5)$$
    $$ = e^f + e^{-f} (6)$$
    $$ {d \over d \gamma}\Gamma^{-1} = e^{f} - e^{-f} (7)$$
    (1) is my son in acasicrecord with his idea, (2),(3) are myself equation.
    (4),(5) are my son and me with tagge of ideas equation.
    (6),(7) are Takashima Aya and me, and my son in acasicrecode with tolio
    equations.
    These references from Grisha professor and Takeuchi Kaoru.
    I read with this references being hint from these professors.
 
  $$ \log x|_{g_{ij}}^{\nabla L} = f^{f^{'}}, F^f|_{g_{ij}}^{\nabla L}
  |: x \to y, x^p \to y $$ 
  $$  f(x) = \log x = p \log x, f(y) = p \log x  $$
  大域的偏微分方程式と大域的多重積分は、それぞれ次のように成り立っている。
  $$ {d^2 \over df^2}F = F^{f^{'}}\cdot f^{f^{''}},  $$
  $$ \int \int F dx_m = F^f \cdot F^{(f)^{'}} $$
  $$ {d \over df dg}(f,g) = (f \cdot g)^{f^{'} + g^{'}} $$
  大域的部分積分も、次のように成り立っている。
  $$ (F^f \cdot G^g) = \int (f \cdot g)^{f^{'} + g^{'}} $$
  $$ \int \int F\cdot G dx_m = [F^f \cdot G^g] - 
  \int (f \cdot g)^{f^{'} + g^{'}} $$
  大域的部分積分の計算は、
  $$ \int {d \over {df dg}}FG = \int F^{f^{'}}G + \int F G^{g^{'}} $$
  $$ \int F^{f^{'}}Gdx_m = [F^f G^g] - \int F G^{g^{'}}dx_m$$
  大域的商代数の計算は、
  $$ \left({F(x) \over G(x)}\right)^{(fg)^{'}} = {F^{f^{'}}G - FG^{g^{'}}
  \over G^{g}} $$
  大域的偏微分方程式は、縮約記号を使うと、
  $$ {d \over df dg}FG = {d \over df_m}FG $$
  $$ {\partial \over \partial f_m}FG = F^{f^{\mu\nu}}\cdot G
  + F \cdot G^{g^{\mu\nu}}, \int F dx_m = F^f $$
  多様体による大域的微分と大域的積分が、エントロピー式で統一的に表せられる。

  $$ 1 + \tan^2 x = {1 \over {\cos^2 x}}$$
  $$ \int\int{1 \over {(x\log x)}^{2}}dx_m
  = \int\int{(1 + \tan^2 x)}dx_m = {d \over {d f}}{F}=F^{f^{'}} $$ 
  $${d \over {d f}}{(\cos^2 x + \sin^2 x + 
  {\left({{\sin x} \over {\cos x}}\right)}^2)} = \pi + e^{f} + e^{-f} $$
  This proof of generate of deprivate and integral calcurate is equal with
  computed of faint in deprivate and integral manifold.
  And, this proof compute is under of explain,

 Rotate with pole of dimension is converted from other sequence of element,  
this atmosphere of vacume from being emerged with quarks of being constructed
with Higgs field, and this created from energy of zero dimension are begun with
universe and other dimension started from darkmatter that represented with 
big-ban.
Therefore, this element of circumstance have with quarks of twelve lake with
atoms.

  $$ (dx, \partial x)\cdot (\epsilon x, \delta x) = \begin{pmatrix}
	  1 & 0 \\
  0 & - 1 \end{pmatrix}^{1 \over 2}  $$
  $$ {d \over df} F = m(x) $$

 And this phenomounen is super symmetry theorem built with quarks of element, 
also this chemistry of mechanism estimate from physics of operatsion.
These response of mechanism chain of geometry into space being emerged with
creature and univse of existing of combination.

  This converted with dimension of element also emerged from imaginary and 
reality of pole's space.

  And this pole of transport of dimension belong with vector of time has 
with one rout of sequecence. Quanum physics also belong with other vector 
of time has with imaginary rout of sequecense.

 This sequence of being estimated with non fluer of time, and this space 
of element have with gravity and antigravity of power.
Other vecor of time is antigravity rout of sequecense.

 Weak electric theory is estimate from time has with one rout of sequecense,
this topology of chain is resulted from time of rout ways.

     $$ \square ({{\sigma_1 + \sigma_2} \over 2 }) = [3 \pi (\chi,x) \circ f(x),
   \sigma(x)] $$
   $$ \square \psi = 8 \pi G T^{\mu\nu} $$
   $$ {d \over dt}g_{ij}(t) = - 2 R_{ij} $$
   $$ \nabla \psi^2 = 4 \pi G \rho $$
   $$ \square (\sigma_1 + \sigma_2) = [6\pi(\chi,x)\circ f(x), \sigma(x)] $$
   $$ = [i \pi(\chi,x),f(x)] $$
   $$ \square \psi = [12 \pi (\chi,x)\circ f(x), \sigma(x)] $$
   $$ {d \over df} F = \int e^{-f}[- \Delta v + R_{ij} v_{ij}+ 
   \nabla_i \nabla_j f  + v \nabla_i \nabla_j + 2<f,h> + (R + \nabla f)
   ({v \over 2} - h) ] $$
   $$ =[i \pi(\chi,x), f(x)] $$

   These equation is reminded time pass rout of one rout way of forms,
   and this rout of time ways which go for system from future and past.
   Therefore, this resulted system of time mechanism is one true flow
   that weak electric theorem oneselves.
   and moreover, one rout time way of forms is reverse with antigravity of
   time system. This also spectrum focus is true that Maxwell theorem and
   strong boson unite with antigravity, this unite is essense on the contrary
   from weak electric theorem, this theorem called for time rout forms is
   strong electric theorem.
   This two theorem is united with quantum physics that no time flow system.

   $$ f^{-1}(x)xf(x) = 1,H_m = E_m \times K_m $$


   The non-commutative theorem is constructed from world line surface that
   this complex manifold estimate with rolanz atractor, and this string
   theorem have with one world of universe mate six quarks and other world of
   dimension mate other element of six quarks.
   These particle quarks is built with super symmetry space of dimension.

   $$ i = (1,0)\cdot (1,0), e^{i\theta} = \cos \theta + i \sin \theta $$
   $$ e^{-i \theta} = \cos \theta - i\sin \theta $$
   $$ \sin \theta = {e^{i \theta} - e^{-i \theta} \over 2i} $$
   $$ \sin i \theta = {e^{-\theta} - e^{\theta} \over 2} $$
   $$ \pi (\chi,x) = \cos \theta + i \sin \theta $$

   In this equations, two dimension redestructed into three dimension,
   this destroy of reconstructed way is append with fifth dimension.
   This deconstructed way of redestructe is arround of universe attached
   with three dimension, this over cover call into fifth dimension.

   $$ R(-\alpha) = \begin{pmatrix} \cos \alpha & \sin \alpha \\
   - \sin \alpha & \cos \alpha \end{pmatrix} $$
   $$ R(\alpha) = \begin{pmatrix} \cos \alpha & - \sin \alpha \\
   \sin \alpha & \cos \alpha \end{pmatrix} $$
   $$ R(\alpha)M R(-\alpha) = 
   \begin{pmatrix} \cos \alpha & - \sin \alpha \\
   \sin \alpha & \cos \alpha \end{pmatrix} \begin{pmatrix} 1 & 0 \\
   0 & - 1 \end{pmatrix} \begin{pmatrix} \cos \alpha & \sin \alpha \\
   - \sin \alpha & \cos \alpha \end{pmatrix} $$
   $$ = \begin{pmatrix} \cos^2 \alpha - \sin^2 \alpha & 2 \sin \alpha \cos 
	   \alpha \\
	   2 \sin \alpha \cos \alpha & - \cos^2 \alpha + \sin^2 \alpha
   \end{pmatrix} $$
   $$ = \begin{pmatrix} \cos 2 \alpha & \sin 2 \alpha \\
   \sin 2 \alpha & - \cos 2 \alpha \end{pmatrix} $$

   $$ \begin{pmatrix} \cos \alpha & - 1 \\
	   1 & - \sin \alpha \end{pmatrix} \begin{pmatrix} \cos \alpha \\ 
		   \sin \alpha
   \end{pmatrix} = \begin{pmatrix} 1 & 0 \\
   0 & - 1 \end{pmatrix} $$
   $$ \begin{pmatrix} \cos \alpha & \sin \alpha \\
   \sin \alpha & - \cos \alpha \end{pmatrix} \begin{pmatrix}
   x \\ y \end{pmatrix} = \begin{pmatrix} 1 & 0 \\
   0 & - 1 \end{pmatrix} $$
   $$ \lim_{\theta \to 0} \begin{pmatrix} \sin \theta \\ \cos \theta 
   \end{pmatrix} \begin{pmatrix} \theta & 1 \\
   1 & \theta \end{pmatrix} \begin{pmatrix} \cos \theta \\ \sin \theta
	   \end{pmatrix} = \begin{pmatrix} 1 & 0 \\ 0 & - 1 \end{pmatrix} $$
		   $$ f^{-1}xf(x) = 1 $$
		   $$ (\log x)^{'} = {1 \over x}, x^n + y^n = z^n  $$
		   $$ x^n = - y^n + c, nx^{n-1} = -n y^{n-1}y^{'} $$
		   $$ y^{'} = {{n x^{n-1}} \over {n y^{n-1}}} $$
		   $$ = {x^{n-1} \over y^{n-1}}= - {y \over x}\cdot
		   ({x \over y})^n $$
		   $$ - {\cos x \over (\cos x)^{'}}(\sin x)^{'} = z_n $$
		   $$ z^n = - 2 e^{x \log x} $$
		   $$ \lim_{x \to \infty} f(x) = a, \lim_{y \to \infty}
		   f(y) = b, \lim_{x,y \to \infty}\{f(x) + f(y)\} = a + b  $$
		   $$ \lim_{z \to 0}f(z) = c, \delta \int z^n = {d \over dV}
		   x^3 $$
		   $$ \lim_{x \to \infty} = c - \lim_{y \to \infty}f(y)
		   \lim_{x \to \infty} f(x) + \lim_{y \to \infty}f(y) = 
		   \lim_{z \to \infty}f(z) $$
		   $$ z^n = \cos n \theta + i \sin n \theta $$
		   $$ = - 2 e^{x \log x} $$

  These equation is gravity and antigravity equation. 

	   $$ {d \over d\sigma}\begin{bmatrix}{(\sigma_1 + \sigma_2) 
	   \over 2}\end{bmatrix}  $$
	   $$ = \sigma(\upharpoonleft \downharpoonleft) + \sigma(
	   \upuparrows)+ \sigma(\downdownarrows)+ 
	   \sigma(\leftleftarrows)
	   + \sigma(\rightrightarrows) $$
	   $$ {d \over df}{\int \int {1 \over (x \log x)^2}}dx_m
	   = \sigma(\upharpoonright) $$
	   $$ {d \over df}{\int \int {1 \over (y \log y)^{1 \over2}}}dy_m
	   = \sigma(\downharpoonright) $$
	   $$ \sigma(\leftarrowtail) + \sigma(\rightarrowtail)
	   = \int e^{-f}[-\Delta v + R_{ij}v_{ij} + \nabla_{i}\nabla_{j}v
	   + v\nabla_{i}\nabla_{j} + 2 <f,h> + (R + \nabla f)(v - 
	   {h \over 2})]
	   $$
	   $$ \sigma(\downharpoonleft) = \sigma(\upharpoonright 
	   \downharpoonright + \upuparrows 
	   + \rightrightarrows)$$
	   $$ \sigma(\upharpoonleft) =  \sigma(\upharpoonright 
	   \downharpoonright + \downdownarrows 
	   + \leftleftarrows) $$
         
	   $$ \text{weak electric theorem} = \sigma(\upharpoonleft) $$
	   $$ \text{strong electric theorem} = \sigma(\downharpoonleft) $$

	   These equation is represented with topology of string model,
	   and weak electric theorem is constructed with Maxwell theorem 
	   and weak boson, gravity that estimate with this three power united.
	   Moreover, strong boson and Maxwell theorem, antigravity that
	   also estimate with this three power united.
	   This two united power is integrated from gravity and antigravity.
	   Then this united power is zeta function.


      Non certain theorem is component with quamtum computer, this system
   emelite with contempolite of eternal universe asterise and
   secret of database.

   不確定性理論の、宇宙の準同型写像の、部分群と全射としての、
   閉３次元多様体のエントロピー値としての、暗号としての
   量子コンピューターが、Jones多項式を形成しているオイラーの公式
   の虚数として、ダミーを入れると、　Rsa暗号の公開鍵として、
   成り立っている量子暗号が、この量子コンピューターと一緒になって、
   ゼータ関数と量子群を大域的微分方程式を、これらの方程式から
   統合されて、大域的トポロジーが、数学の数論と幾何学、解析学として、
   物理学と言語学、情報科学、すべての理論に統合されて、
   ゼータ関数が、統一場理論の発生する式になっている。
   この理論を彩さんとグリーシャさん、ナッシュさん、トビーケリーさん、
   小方くん、益川先生、南部先生、岡本教授、リサ・ランドール教授、
   竹内先生、私の子ども、今までの人たちが、集まって、この理論ができている。
   量子的な微分・積分が、きっかけにもなっている。
   アイデアと計算方法は、情報空間の子どもと、お姉さんとお父さんと
   先生たちで、統一場理論としては、グリーシャさんと竹内薫先生と
   彩さん、岡本教授がきっかけになっている。


      This information of quantum computer ansterise in secret of each data,
   and this system is entersteam by Jones manifold of equation,
   more also this intersect with system of pair of universe and the other
   dimension with Jones manifold of equation pawn for non access of 
   chain in diseable of database.
      And this insterm of system built with 
   $$ e^{- \theta} = e^f + e^{-f}, e^{i \theta} = \cos \theta +
   i \sin \theta $$
   $$ {d \over df}F = e^f + e^{-f} = 2 i \sin (i x \log x) $$
   These equations are atom of dense with norm and energe of non certain
   theorem, and this system component with also Jones manifold of equation.
   This equation addesent with a dummy data in non integrate of relativity
   theory for quantum physics and access to varnished with aseterise for
   important of descover with secret of data.
   This dummy of information of quantum secret datas are non commutative
   equation of declinate anterave and distarbration 
   into universe and other dimension
   for accesslable with resterbled with quantum equation of 
   deep of Riemann metrics, more also this system is computed with 
   after all intersected with eternal data into Zeta function of
   global differential equation and Higgs fields with time system enterprises.
   Global of open key in Rsa secret system intersect with this base of
   this technolgie with all of data for established with information 
   technology and quamtum computer created with dummy of integral 
   non relativity of quantum equation, this focus is 
   diserble with quantum equation of a data for dummy into 
   Rsa secret data of non commutative equation, this system
   comute with after all established with information of quantum physics.

   $$ e^{i \theta} = \cos \theta + i \sin \theta $$
   this equation is based into Jones manifold equation.
   $$ e^{- \theta} = e^f + e^{-f}, {d \over df}F + \int C dx_m 
   = 2 ( \cos ( i x \log x) + i \sin (i  x \log x)) $$
   These equation are based with reco level theory and field of physics,
   and these system are 
   $$ ||ds^2|| = \int \mathrm{exp}[{- L (x) \over \sqrt{2 q \tau}}]dx_m
   + \mathcal{O}(N^{-1}) $$

   $ \mathcal{O}(N^{-1}) $ is dummy data of intersected of system.

   These system are chain with non certain theorem component with
   aspe experiement mechanism clearlity of quantum teleportation physics 
   and this system is pair of existanse with Jones manifold equation
   and imarginary circle equation more also these equations are 
   rhysmiary equation excluded with zeta function and quantum equation,
   and this system commutave with each attitude of pair in Jones manifold
   equation.

   These theorem combuild with Lie algebra and catastorophe  
   of summative group in eternal Garois theory,
   and this reasons of between Poancare conjecture in zero dimension 
   and eternal group with nesseciry and mourdigate of conditions.


   $ AdS_5 $ equation are built with cercumention dimension of assemble
   D-brane and information of quantum physics with secret system,
   and this system construct with a certain theorem.

   $$ ||ds^2|| = \eta_{\mu\nu} + \bar{h}(x) $$ 
   This dimension esterned with equals of dummy data in the other dimension,
   and this value of data declined into only universe,
   then this sertation of universe is possibilty equation.
   After all this anserration of univese addsented with the other dimension
   escourt for certain theorem and all of addsented equation
   are three manifold entropy equation,
   this equation constructed wih zeta function and quantum equation.

   $$ ||ds^2|| = 8 \pi G \left({p \over c^3} + {V \over S}\right) $$
   $$ 2 \sqrt{V \over S}= {d \over df}\int \int{1 \over (x \log x)^2}dx_m $$
   $$ {d \over df} \int \int{ 1 \over (x \log x)^2}dx_m = \int dx_m + x^2 $$
   $$ \ge \int dx_m $$
   $$ {d \over df} \int \int{1 \over (x \log x)^2}dx_m \ge {1 \over 2}i $$

     Quantum computer concluded with component of being describe with 
     estimate for each atom conditions. This pair of condition is
     entertained for orbital of circle with Euler equation.
     This pair of orbital area estertate with two of circumentations.
     These system concerned with hyper synchronized of all of possibility
     resulted for being computing of all of area in Euler equation.
     This cercuit with equations are that information of quantum physics
     esterned with Rsa secret datas system.


     These system of stem are constructed with declinate anterave of
     Lambda driver and possibilty computer for resolved with synchronize
     of distarbrate non certain theorem, and this concerned with
     system component of Euler equation.
     This estimate of orbital construct with curbit of quarters.
     Artificial intelligence is concluded with Euler equation.
     Zeta function also integrate with this equation.
     これらの理論の中枢の,中核となっている、オイラーの公式の群とも
     なっている、過冷却金属の原理の、高熱を急激に冷ます、電子を
     集める、ラムダ・ドライバーにもなっている、不確定性理論の原理
     にもなっている、 暗号を解読するのを防ぐのに、
     クオータのキュービットの、遷移元素を作り出す
     原子の仕組みにもなっている、この理論が、量子コンピュータと
     人工知能の中枢となっていて、オイラーの公式の虚数ともなっている、
     このJones多項式の統合にゼータ関数が使われている。
 

 Hilbert manifold equals with Von Neumann manifold, and this fields is concluded with Glassman manifold,
the manifold is dualty of twister into created of surface built.
These fields is used for relativity theorem and quantum physics.
$$ ||ds^2|| = \lim_{x \to \infty} [\delta(x) \int \int \int \pi \left( \sum_{k=0}^{\infty}
		{{}^n\sqrt{p},x \over n} \right)
^{1 \over 2}d \tau]^{\mu\nu} $$
  This norm is component of fields in Hilbert manifold of space theorem rebuilt with
Mebius space in Gamma Function on Beta function fill into power of rout fundamental
group.
And this result with $AdS_5$ space time in Quantum caos of Minkowsky of manifold in
abel manifold.
Gravity of metric in non-commutative equation is Global differential equation
conbuilt with Kaluza-Klein space. Therefore this mechanism is $T^{\mu\nu}$ tensor
is equal with $R^{\mu\nu}$ tensor. And This moreover inspect with laplace operator
in stimulate with sign of differential operator.
Minus of zone in Add position of manifold is Volume of laplace equation rebuilt with
Gamma function equal with summative of manifold in Global differential equation,
this result with setminus of zone of add summative of manifold, and construct with
locality theorem straight with fundamental group in world line of surface, this
power is boson and fermison of cone in hyper function.

$$ V(\tau) = [f(x),g(x)] \times [f^{-1}(x), h(x)] $$
$$ \Gamma(p,q) = \int e^{-x} x^{1 - t}dx $$
$$ = \beta(p,q) $$
$$ = \pi (f(\chi,x),x) $$
$$ ||ds^2|| = \mathcal{O}(x)[(f(x) \circ g(x))^{\mu\nu}]dx^{\mu}dx^{\nu} $$
$$ = \lim_{x \to \infty} \sum_{k=0}^{\infty}a_k f^k $$
$$ G^{\mu\nu} = {\partial \over \partial f}\int [f(x)^{\mu\nu} \circ G(x)^{\mu\nu}dx^{\mu}dx^{\nu}]^{\mu\nu}dm $$
$$ = g_{\mu\nu}(x)dx^{\mu}dx^{\nu} - f(x)^{\mu\nu}dx^{\mu}dx^{\nu} $$
$$ [i \pi(\chi,x),f(x)] = i \pi f(x) - f(x) \pi(\chi,x) $$
$$ T^{\mu\nu} = (\lim_{x \to \infty} \sum_{k=0}^{\infty}\int \int [V(\tau) \circ S^{\mu\nu}(\chi,x)]dm)
^{\mu\nu}dx^{\mu}dx^{\nu} $$
$$ G^{\mu\nu} = R^{\mu\nu}T^{\mu\nu} $$
$$ \begin{vmatrix} D^m & dx \\
dx & \sigma^m \end{vmatrix} \begin{vmatrix} \cos \theta & - \sin \theta \\
\sin \theta & \cos \theta \end{vmatrix} \begin{vmatrix} x \\
y \end{vmatrix} = \begin{pmatrix} 1 & 0 \\ 
0 & - 1 \end{pmatrix}^{1 \over 2} $$
$$ \sigma^m \begin{bmatrix} \delta(x) & - 1 \\
1 & \epsilon(x) \end{bmatrix}^{1 \over 2} = \begin{pmatrix} i & 0 \\
0 & - i \end{pmatrix} $$
$$ V(M) = {\partial \over \partial f} ({}^N \int [f \smallsetminus M]^{\oplus N})^{\mu\nu} dx^{\mu}dx^{\nu} $$
$$ V(M) = \pi(2 \int \sin^2 dx)\oplus {d \over df}F^M dx_m $$
$$ \lim_{x \to \infty}\sum_{k=0}^{\infty} a_k f^k = \int (F(V)dx_m)^{\mu\nu}dx^{\mu}dx^{\nu} $$
$$ \bigoplus_{k=0}^{\infty}[ f \smallsetminus g] = \vee(M \wedge N) $$
$$ \pi_1 (M) = e^{-f 2 \int \sin^2 x dm} + O(N^{-1}) $$
$$ = [i \pi(\chi,x),f(x)] $$
$$ M \circ f(x) = e^{-f \int \sin x \cos x dx_m} + \log (O(N^{-1})) $$

 Non-Symmetry space time.
 $$ {d \over dL}V(\tau) = {d \over df}\int \int_M {1 \over (x \log x)^2}dx_m + {d \over df}\int \int_M
 {1 \over (y \log y)^{1 \over 2}}dy_m $$
 $$ \epsilon S(\nu) = \square_v \cdot {\partial \over \partial \chi}({}^5 \sqrt{\wedge g^2})d\chi $$
 Differential Volume in $AdS_5$ graviton of fundamental rout of group.
 $$ \wedge (F^m_t)^{''} = {1 \over 12}g_{ij}^2$$
 Quarks of other dimension.
 $$ \pi(V_{\tau}) = {e^{-(\sqrt{\pi \over 16}\log x)}}^{\delta} \times {1 \over (x \log x)} $$
 Universe of rout, Volume in expanding space time.
 $$ {d \over dt}(g_{ij})^2 = {1 \over 24}(F^m_t)^2 $$
 $$ m^2 = 2 \pi T \left({26 - D_n \over 24}\right) $$
 This quarks of mass in relativity theorem, and fourth of universe in three manifold of one dimension surface,
 and also this integrate of dimension in conbuilt of quarks.
 $$ g_{ij} \wedge \pi(\nu_{\tau}) = e^{-2 \pi T |\psi|}[\eta_{\mu\nu} + \bar{h}_{\mu\nu}(x)]dx^{\mu}dx^{\nu} 
 + T^2 d \psi^2 $$
 Out of rout in $AdS_5$ space time.
 $$ ||ds^2|| = g_{ij}\wedge \pi(\nu_v) $$
 $AdS_5$ norm is fourth of universe of power in three manifold out of rout.


\section{Entropy}
   Already discover with this entropy of manifold is blackhole of energy, these entropy is mass of energy has
norm parameter. And string theorem $ T_{\mu\nu}$  belong to $ \rho $ energy, this equation is same Kaluze-Klein theorem.


 $$ \nabla \phi^2 = 8 \pi G({p \over c^3} + {V \over S})$$

 $$  y = x $$
 $$  y = (\nabla \phi)^{1 \over 2} $$

  $$ {\Delta E \over {1 \over 2\sqrt{2 \pi G}}} = {p \over c^3} + \rho $$

  $$ ds^2 = g_{\mu\nu}(x)dx^{\mu}dx^{\nu} + \phi^2(x)\kappa^2 A_{\mu\nu}(x)dx^{\mu})^2 $$
  $$ ds = (g_{\mu\nu}(x)dx^{\mu}dx^{\nu} + \phi^2(x)(\kappa^2 A_{\mu\nu}(x)dx^{\mu})^2)^{1 \over 2} $$

  $$ {\Delta E \over {1 \over 2\sqrt{2 \pi G}} - \rho} = c^3, {p \over 2 \pi} = c^3 $$
    

  $$ ds^2 = e^{-2kT(x)|\phi|}[\eta_{\mu\nu} + \bar{h}_{\mu\nu}(x)]dx^{\mu}dx^{\nu} + T^2(x) d\phi^2 $$

  $$ f_{z} = \int \begin{bmatrix} \sqrt{\begin{pmatrix} x_{1} & x_{2} & x_{3} \\
     y_{1} & y_{2} & y_{3} \end{pmatrix}
     \circ
   \begin{pmatrix} x_{1} & x_{2} & x_{3} \\
     y_{1} & y_{2} & y_{3} \end{pmatrix}}_{}
     \end{bmatrix}dxdydz $$

     $$ {\partial \over \partial f} \int (\sin2x)^2dx = ||x - y||^2 $$
	$$ ||\int [\nabla_{i}\nabla_{j}f]dm||^{{1 \over 2} + iy} = \mathrm{rot}(\mathrm{div},E,E_1) $$
	Maxwell of equation in fourth of power.
	$$ = 2 <f,h>, {V(x) \over f(x)} = \rho (x) $$
	$$ \int_M \rho (x)dx = \square \psi, - 2 <g, h> = \mathrm{div}(\mathrm{rot} E,E_1) $$
	$$ = - 2 R_{ij} $$
	Higgs field of space quality.
	$$ \mathcal{O}(x) = ||{\nabla_{i}\nabla_{j} \over S^2}\int [\nabla_{i}\nabla_{j}f\cdot g(x)dxdy]d\psi|| $$
	$$ = \int (\delta(x))^{2 \sin \theta \cos \theta} \log \sin \theta d \theta d\psi $$
	Helmander manifold of duality differential operator.
	$$ \delta \cdot \mathcal{O}(x) = [{||\nabla_{i}\nabla_{j} f d \eta|| \over 
	\int e^{2 \sin \theta \cos \theta} \cdot \log \sin \theta d \theta}] $$	
	Open set group in subset theorem, and helmandor manifold in singularity.
	$$ i\hbar \psi = ||\int {\partial \over \partial z} [{i(xy + \overline{yx}) \over z - \bar{z}}]dmd\psi|| $$
	Euler number. Complex curvature in Gamma function. Norm space.
	Wave of quantum equation in complex differential variable into integrate of rout.
	$$ \delta(x) \int_C R^{2 \sin \theta \cos \theta} = \Gamma(p + q) $$
	$$ (\square + m) \cdot \psi = (\nabla_{i}\nabla_{j}f|_{g_{\mu\nu(x)}} + v \nabla_{i}\nabla_{j}) $$
	$$ \int[m(x)(\mathrm{rot}\cdot\mathrm{div}(E,E_1))]dmd\psi $$
	Weak electric power and Maxwell equation combine with strong power.
	Private energy.
	$$ G_{\mu\nu} = \square \int \int \int (x,y,z)^3dxdydz $$
	$$ = {8 \pi G \over c^4}T^{\mu\nu} $$
	$$ {d \over dV}F = \delta(x) \int \nabla_{i}\nabla_{j}f d \eta_{\mu\nu} $$
	Dalanvelsian in duality of differential operator and global integrate variable operator.

	$$ x^n + y^n = z^n, \delta (x) \int z^n = {d \over dV}z^3, (x,y)\cdot(\delta^m,\partial^m) $$
	$$ = (x,y)\cdot(z^n,f) $$
	$$ n \bot x, n \bot y $$
	$$ = 0 $$
	Singularity of constance theorem.

	$$ \vee(\nabla_{i}\nabla_{j}f) \cdot \mathrm{XOR}(\square \psi)= {d \over df}\int_M F dV $$
Non-cognitive of summuate of equation and add manifold create with open set group theorem.
This group composite with prime of field theorem. Moreover this equations involved with symmetry dimension created.
	$$ \partial (x) \int z^3 = {d \over dV}z^3, \sum_{k=0}^{\infty}{1 \over (n + 1)^s} = \mathcal{O}(x) $$
	Singularity theorem and fermer equation concluded by this formula.
	$$ \int \mathcal{O}(x)dx = \delta(x)\pi(x)f(x) $$
	$$ \mathcal{O}(x) = \int \sigma(x)^{2\sin \theta \cos \theta}\log \sin \theta d\theta d\psi $$
	$$ y = x $$
	$$ \mathcal{O}(x) = ||{\nabla_{i}\nabla_{j}f \over S^2} \int [\nabla_{i}\nabla_{j}f \circ g(x)dxdy]|| $$
	$$ \lim\limits_{n \to 1}\sum_{k=0}^{\infty}{a_k \over a_{k+1}} = (\log \sin \theta dx)^{'} $$
	$$ = {\cos \theta \over \sin \theta} $$
	$$ = {x \over y} $$
	$$ \lim\limits_{n \to 1}\sum_{k=0}^{\infty}a_k f^k = {1 \over 1 - z} $$
	Duality of differential summuate with this gravity theorem.
	And Dalanverle equation involved with this equation.
	$$ \square \psi = {8\pi G \over c^4}T^{\mu\nu} $$
	$$ {\partial \over \partial f}\square \psi = 4 \pi G \rho $$
	$$ \int \rho (x) = \square \psi, {V(x) \over f(x)} = \rho(x) $$
	Dense of summuate stimulate with gravity theorem.

  今、大域的積分多様体に微分を入れると、大域的部分積分多様体が
   できる。ガンマ関数の大域的微分多様体と同型らしい。　


      大域的多重積分と大域的無限微分多様体を定義すると、
   Global assemble manifold defined with infinity of differential
   fields to determine of definition create from Euler product and
   Gamma function for Beta function being potten from integral manifold
   system. Therefore, this defined from global assemble manifold from
   being Gamma of global equation in topological expalanations.

   $$ \int f(x)dx = \int \Gamma(\gamma)^{'}dx_m $$
   $$ = 2 (\cos (i x \log x) - i \sin(i x \log x)) $$
   $$ \left({{\int f(x) dx} \over \log x}\right) = \lim_{\theta \to \infty}
    \left({{\int f(x) dx} \over \theta}\right) = 0,1 $$
    $$ e^{i \theta} = \cos \theta + i \sin \theta $$
    $$ \left({\int f(x) dx}\right)^{'} = 2 (i \sin(i x \log x) - \cos (i x \log
    x)) $$
    $$ = 2(- \cos(i x \log x) + i \sin(i x \log x)) $$
    $$ (\cos (i x \log x) - i \sin(i x \log x))^{'} $$
    $$ = {d \over {d {e^{i\theta}}}}\left((\cos, - \sin )\cdot
    (\sin, \cos)\right)  $$
    $$ \int \Gamma(\gamma)^{''}dx_m = \left({\int \Gamma(\gamma)^{'}}dx_m
    \right)^{\nabla L} = \left({{\int \Gamma dx_m} \cdot {{d \over d\gamma}
    \Gamma}}\right)^{\nabla L} \le 
    \left({\int \Gamma dx_m + {d \over d\gamma}\Gamma}\right)^{\nabla L} \le 
    \left({e^{f} - e^{-f}} \le {e^{-f} + e^{f}}\right)^{'}  $$
    $$ = 0,1 $$

    Then, the defined of global integral and differential manifold
    from being assembled world lines and partial equation of deprivate
    formula in Homology manifold and gamma function of global topology
    system, this system defined with Shwaltsshild cicle of norm from
    black hole of entropy exsented with result of monotonicity for
    constant of equations.
    以上 より、大域的微分多様体を大域的 2重微分多様体として、
    処理すると、ホモロジー多様体では、種数が１であり、特異点では、種数が０
    と計算されることになる。ガンマ関数の大域的微分多様体では、
    シュバルツシルト半径として計算されうるが、これを大域的2重微分で
    処理すると、ブラックホールの特異点としての解が無になる。




	Abel拡大 $ K/k $に対して、
	$$ f = {\pi}_{p}{f_p} $$
        類体論
	Artin記号を用いて、
       $$ \left({{\alpha, K/k} \over {p}}\right) 
       = \left({{K/k} \over {b}}\right) (\in G) $$
       $ {\alpha/ \alpha_0} \equiv 1 \pmod {f_p} $, 
      $ \alpha_0 \equiv 1 \pmod {f f_p^{-1}} \to  
      \alpha \in k $
      $ (\alpha_0) = {p}^{\alpha}b, \text{
	      $p$ と$b$は互いに素} $
      $b \to \text{相対判別式$\delta K/k$で互いに素}$
      この値は、補助数${\alpha_0}$の値の取り方によらずに、
      一意的に定まる。
      $$ \left({{\alpha, K/k} \over {p_{\infty}^{(j)}}}\right)  = 
      \text{1または0} $$

      これらをまとめた式が、Hilbertの剰余記号の判別式
      $$ \mathcal{\pi}_{p}\left({{\alpha,b} \over 
      {p}}\right) = 1 $$
      であり、
    　この式たちから、代数幾何の種数のノルム記号である、
$$ ||ds^2|| = \lim_{x \to \infty} [\delta(x) \int \int \int \pi \left( \sum_{k=0}^{\infty}
		{{}^n\sqrt{p},x \over n} \right)
^{1 \over 2}d \tau]^{\mu\nu} $$
     が求まり、
     $$ p^{\alpha} n = {}^n \sqrt{p} $$
     $$ n^{{}^n\sqrt{p}} = \bigoplus{(i \hbar^{{{\nabla)}}^{\oplus 
     L}}} $$
     $$ = n^{p^{{1 \over n}}} = n^{{-n}^{p}}  $$
     $$ = \int \Gamma(\gamma)^{'}dx_m = e^{-x \log x} $$
     となり、
     kの素イデアルの密度Mに対して、
     $$ \lim_{s \to {1 + 0}} \sum_{{p \in M}} {
	     1 \over (N(p))^{s}}/ \log {1 \over 
     {s - 1} } $$
     $$ = \text{Mの密度(density)} $$

     $$ \alpha(f{d \over dt},g{d \over dt}) = \int_s \begin{vmatrix}
	     f^{'} & f^{''} \\
	     g^{'} & g^{''} \end{vmatrix}dt,
	     \mathcal{B}(f{d \over dt},g{d \over dt},h{d \over dt}) =
	     \int_s \begin{vmatrix} f & f^{'} & f^{''} \\
		     g & g^{'} & g^{''} \\
	     h & h^{'} & h^{''} \end{vmatrix}dt $$
    これらは、Gul'faid-Fuksコホモロジーの概念を局所化することにより、
    形式的ベクトル和のつくるコホモロジーとして、導かれている。
$$ || \int f(x) ||^2 \to \int \pi r^2 dx \cong V(\tau) $$

 This equation is part of summulate on entropy formula.

 $$ {\int |f(x)|^2 \over f(x)} dx $$
$$V(\tau) \to mesh $$

 $$ \int \begin{vmatrix} x_1 & x_2 & x_3 \\
    y_1 & y_2 & y_3 \\
    z_1 & z_2 & z_3 \end{vmatrix} dv({\tau}) $$
 $$ \begin{vmatrix} x_1 & x_2 & x_3 \\
    y_1 & y_2 & y_3 \\
    z_1 & z_2 & z_3 \end{vmatrix}_{dx=v}  $$


   $$ z_y = a_x + b_y + c_z $$
$$ dz_y = d(z_y) $$

$$ [f, f^{-1}] = ff^{-1} - f^{-1}f $$

  Universe of extend space has non-integrate rout entropy in sufficient of mass in darkmatter,
this reason belong for cohomology and heat equation concerned. This expanded mass match is darkmatter
of extend. And these heat area consume is darkmatter begun to start big-ban system invite of universe.

$$ V(\tau) = \int \tau(q)^{-{n \over 2}} \mathrm{exp}(-{1 \over \sqrt{2 \tau(q)}}L(x)dx) + O(N^{-1}) $$
$$ {1 \over \tau}({N \over 2} + \tau(2 \Delta f - |\nabla f|^2 + R) + f)\mathrm{mod}N^{-1} $$
	$$ \Delta E = -2(T - t)|R_{ij} + \nabla_{i} \nabla_{j} f - {1 \over 2(T - t)}g_{ij}|^2 $$
	$$ {d \over df}F = {2\int (R + \nabla_{i} \nabla_{j} f)^2 \over - (R + \Delta f)}dm $$
	$$ {d \over dt}g_{ij}(t) = - 2 R_{ij} $$
	$$ dx = (g_{\mu\nu}(x)^2dx^2 - g_{\mu\nu}(x)dxg_{\mu\nu})^{1 \over 2} $$
	$$ ds^2 = - N(r)^2dt^2 + \psi^2(r)(dr^2 + r^2d\theta^2) $$
  $$ f_{z} = \int \begin{bmatrix} \sqrt{\begin{pmatrix} x_{1} & x_{2} & x_{3} \\
     y_{1} & y_{2} & y_{3} \end{pmatrix}
     \circ
   \begin{pmatrix} x_{1} & x_{2} & x_{3} \\
     y_{1} & y_{2} & y_{3} \end{pmatrix}}_{}
     \end{bmatrix}dxdydz $$

     $$ \sum^{\infty}_{n=0}a_1x^1 + a_2x^2 \dots a_{n-1}x^{n-1} \to \sum^{\infty}_{n=0} a_nx^n \to \alpha $$

$$ f = n\nu\lambda , \lambda = {x \over l},
 \int dn \nu \lambda = f(x), 
 xf(x) = F(x),
 [f(x)] = \nu h$$


    石岡先生のところで見て、石川先生と河相先生、脇本先生のところで、
    確かめている。

    $$ {2 e^{x_n t} \over {e^{t} + 1}} = \sum_{n=0}^{\infty} E_n(x){{t^x} 
    \over {n!}} $$
    $$ E_n(x) = \sum_{k=0}^{n} \begin{pmatrix} n \\
    k\end{pmatrix}  a_n x^{n -k} $$

	   $$ y =  1 + \tan^2 x = {1 \over {\cos x}} $$
	   $$ y^{'} = {{\cos^{'}x} \over {\cos^2 x}} $$
           $$ -{{\cos x} \over (\cos x)^{'}}(\sin x)^{'} = z_n $$
	   $$ {1 \over y^{'}} = z^2_n$$
           $$ z^n = - 2 e^{x \log x} $$

	   ガンマ関数の大域的積分多様体は、結局は、フェルマーの定理
	   だった。

         代数的計算手法のために$ \oplus L $を使っている。
そのために、冪乗計算と商代数の計算が、乗算で楽に見えるようになっている。
            微分幾何の量子化は、代数幾何の量子化の計算になっている。
  加減乗除が初等幾何であり、大域的微分と積分が、現代数学の代数計算の
  簡易での楽になる計算になっている。
  初等代数の計算は、
  $$ \bigoplus({i \hbar^{\nabla}})^{\oplus L} $$
   $$m \bigoplus({i \hbar^{\nabla}})^{\oplus L} +  
   n\bigoplus({i \hbar^{\nabla}})^{\oplus L} = 
   (m+n) \bigoplus({i \hbar^{\nabla}})^{\oplus L}   $$
   $$ m\bigoplus({i \hbar^{\nabla}})^{\oplus L} -  
   n\bigoplus({i \hbar^{\nabla}})^{\oplus L} = 
   (m-n) \bigoplus({i \hbar^{\nabla}})^{\oplus L}   $$
   $$ \bigoplus({i \hbar^{\nabla}})^{{\oplus L}^m} \times 
      \bigoplus({i \hbar^{\nabla}})^{{\oplus L}^n} =
   \bigoplus({i \hbar^{\nabla}})^{\oplus L^{m+n}} $$
  $$ {{\bigoplus({i \hbar^{\nabla}})^{\oplus L^m}}  \over  
  {\bigoplus({i \hbar^{\nabla}})^{\oplus L^n}}} =  
  \bigoplus({i \hbar^{\nabla}})^{\oplus L^{m\over n}} $$
  大域的計算での微分と積分は、
  $$ \left(\bigoplus({i \hbar^{\nabla}})^{\oplus L}\right)^{df}
  =  {\left(\bigoplus({i \hbar^{\nabla}})^{\oplus L}\right)^{
	  \bigoplus({i \hbar^{\nabla}})^{\oplus L}}}^{'}  $$
  $$ = \bigoplus({i \hbar^{\nabla}})^{\oplus L^{'}} $$
   $$ \int \bigoplus({i \hbar^{\nabla}})^{\oplus L} dx_m $$
   $$ \bigoplus({i \hbar^{\nabla}})^{\oplus L} = n $$
   大域的積分の代数幾何の量子化での計算は、ガンマ関数になっている。
   $$ {n^{L+1} \over {n+1}} = \int e^x x^{1 - t} dx $$
   代数幾何の量子化の因数分解は、加減乗除の式のm,nの
   組み合わせ多様体でのガンマ関数同士の計算になる。
   $$ \left(\bigoplus({i \hbar^{\nabla}})^{\oplus L}+m\right) 
   \left(\bigoplus({i \hbar^{\nabla}})^{\oplus L}+n\right) $$
  $$ = {n^{L+1} \over {n+1}} = t \int (1 - x)^n x^{m - 1} dx $$
  $$ \int x^{m-1}(1 - x)^{n-1} dx $$
  $$ \nabla ({i \hbar^{\nabla}})^{\oplus L}, \bigoplus ({i \hbar^{\nabla}})^{
	   \oplus L}, \square ({i \hbar^{\nabla}})^{\oplus L} $$
   $$ \boxtimes (i \hbar^{\nabla})|_{dx_m}^L, \boxplus (i \hbar^{\nabla}|_
   {dx_m}^L) $$



  $$ x^{{1 \over 2} + iy} = e^{x \log x} $$
  それゆえに、この定数はゼータ関数と微分幾何の量子化を因数にもつ
  素因子分解の式にもなっている。
  Therefore, this product also constructed with differential geometry
  of quantum level equation and zeta function.
  $$ = C $$ 
  そして、この関数はオイラーの定数から広中平祐定理による4重帰納法の
  オイラーの公式からの多様体積分へとのサーストン空間のスペクトラム関数
  ともなっている。
 
  $$ (\psi^{out}(h_1, \dotsb h_n), \psi^{out}(h^{'}_1, \dotsb h^{'}_{{n^{'}}}
  )) $$
  $$ = \delta_{n n^{'}} \sum_{p} \Pi^{n}_{j=1} (h,h_{p}(
  j)) $$
  このHaag-Ruelleの散乱理論からのFucks空間の$ \psi^{out} $の１粒子状態が
  $\sum^{\oplus}R_1(m)$上の$\mathcal{K^{0}}$への閉部分空間への
  $R^{out}$のユニタリ写像として、複素空間の種数の重ね合わせをノルム空間と
  見なして、これが代数幾何の量子化と同型と言えることを、
  これらの式たちから示しているのが、Hilbert空間のノルムである。　


\end{document}
